\section{Verificación y evaluación del modelo}
\label{sec:verificacion-evaluacion}

\subsection{Alcance de las comprobaciones}
La verificación comprueba que el programa resuelva las ecuaciones planteadas. La calibración estima parámetros a partir de datos, como ocurrió en la inversión docente de las probetas. Validar la predicción de la fundación requeriría compararla con temperaturas independientes medidas durante su construcción. Esas observaciones no están disponibles.

\subsection{Comparación con una solución conocida}
Se consideró una barra sin fuentes, de longitud $L=1\,\mathrm m$ y difusividad $\alpha=10^{-4}\,\mathrm{m^2/s}$. Ambos extremos se mantuvieron a $T_b=20\,^\circ\mathrm C$ y se impuso el perfil inicial $T_b+A\sin(\pi x/L)$, con $A=50\,\mathrm K$. Para esas condiciones, la solución exacta de $\partial T/\partial t=\alpha\,\partial^2T/\partial x^2$ conserva la forma del perfil y reduce su amplitud con el tiempo:
\begin{equation}
  T_{\mathrm{ex}}(x,t)=T_b+A\sin\left(\frac{\pi x}{L}\right)
  \exp\left(-\frac{\alpha\pi^2t}{L^2}\right),
  \label{eq:caso-seno-verificacion}
\end{equation}
Con $100$ intervalos, un paso de aproximadamente $0{,}44991\,\mathrm s$ y un horizonte de $1500\,\mathrm s$, el mayor error absoluto entre los nodos y tiempos guardados fue $0{,}0025714\,\mathrm K$. El valor quedó bajo la tolerancia de $0{,}003\,\mathrm K$ del ensayo. También se contrastó la solución superficial de Newton con bisección y se comprobaron las transiciones y el límite de potencia de los circuitos. La barra permite aislar el error de difusión sin confundirlo con el efecto de la geometría de la fundación.

\subsection{Cierre del balance energético}
La energía sensible $E_c$ se calcula sobre las celdas activas respecto de $0\,^\circ\mathrm C$. Su cambio debe recoger tanto el calor generado como la energía que ingresa al colocar hormigón nuevo. Con los aportes acumulados en julios, el residuo es
\begin{equation}
  R_E=E_c-\left(E_{\mathrm{col}}+E_h+E_{\mathrm{aire}}
  +E_{\mathrm{suelo}}+E_{\mathrm{frio}}\right).
  \label{eq:residuo-energia-verificacion}
\end{equation}
$E_{\mathrm{col}}$ corresponde a la colocación y $E_h$ a la hidratación. El término de aire incluye convección y radiación, ya transmitidas a través de las protecciones. Suelo y refrigeración llevan el signo de su aporte al hormigón, de modo que la extracción es negativa. Los flujos internos se cancelan al sumar las celdas.

Se revisaron $144$ balances, dos horizontes para cada una de las $72$ simulaciones finales. El mayor residuo absoluto fue $2{,}48\times10^{-4}\,\mathrm J$. A las $672\,\mathrm h$, el mayor residuo relativo fue $2{,}39\times10^{-15}$ respecto de la suma de magnitudes de los aportes, con denominador mínimo de $1\,\mathrm J$. Esta comprobación respalda la contabilidad energética incluso cuando crece el dominio.

\subsection{Sensibilidad espacial y temporal}
La Tabla~\ref{tab:maximos-verificacion} muestra los extremos de las matrices espaciales. Para evaluar el refinamiento se comparó, además, cada escenario consigo mismo (Tabla~\ref{tab:contrastes-discretizacion}). El mayor cambio de diferencial entre mallas fue $0{,}882\,\mathrm K$ en los $24$ pares. Al reducir solo el paso fue $0{,}00137\,\mathrm K$ en los dos casos estudiados. Ninguna de estas comparaciones cambió la clasificación.
\begin{table}[htbp]
  \centering
  \caption{Máximos del catálogo de la propuesta en las dos mallas}
  \label{tab:maximos-verificacion}
  \small
  \begin{tabular}{ccccc}
    \toprule
    Celda [m] & Paso máximo [s] & Mayor $T_{\max}$ [$^\circ$C] & Mayor $\Delta T_{\max}$ [K] & Excedencias \\
    \midrule
    $0{,}40$ & $300$ & $47{,}4941$ & $16{,}5306$ & $0/24$ \\
    $0{,}20$ & $300$ & $47{,}0322$ & $15{,}7202$ & $0/24$ \\
    \bottomrule
  \end{tabular}
  \captionsetup{skip=2pt}
  \caption*{Fuente: elaboración propia a partir de las simulaciones con ambas mallas.}
\end{table}

\subsection{Márgenes y criterio para aceptar la resolución}
En la malla fina, el diferencial más alto deja $4{,}28\,\mathrm K$ hasta el límite de $20\,\mathrm K$. Ese margen supera el mayor cambio observado por refinamiento espacial. Los dos niveles permiten juzgar la sensibilidad de la decisión, pero son insuficientes para establecer un orden de convergencia o acotar todo el error. Si un refinamiento desplazara un máximo de un lado al otro del límite, habría que refinar la malla o el paso, según el origen de la diferencia, o declarar indeterminado el cumplimiento.

\subsection{Errores e idealizaciones que condicionan la recomendación}
La malla aproxima tanto la posición de las caras como la temperatura a $25\,\mathrm{mm}$. Al mismo tiempo, el suelo a temperatura fija, el contacto perfecto y la activación instantánea de cada capa simplifican procesos que en obra tienen extensión y duración. Una concordancia entre mallas puede coexistir con un sesgo común a ambas.

Las curvas de probeta y los años históricos cubren solo parte de las condiciones posibles. También queda por medir qué potencia pueden extraer los tubos y qué fracción de la radiación intercepta la cubierta instalada. Estas limitaciones orientan las mediciones y las decisiones de retiro que se proponen en la sección~\ref{sec:discusion-control}.
